% concatenated from continblowuparxivfinal.tex
\documentclass[12pt]{amsart}
\usepackage{amssymb}
\usepackage{amsmath,amscd}
\usepackage{amsfonts,latexsym,verbatim,amscd,mathrsfs,color,array}
\usepackage[all,cmtip]{xy}
\usepackage{mathrsfs}
\usepackage{multirow}
\usepackage{graphicx}
\usepackage{amsfonts, amsthm}
\usepackage{bbm}
\usepackage[hmargin=1in,vmargin=1in]{geometry}
\usepackage{mathdots}
\usepackage{stmaryrd}
\usepackage{MnSymbol}
\usepackage{bbm}
\usepackage[hidelinks]{hyperref}

\allowdisplaybreaks

\def\Xint#1{\mathchoice
	{\XXint\displaystyle\textstyle{#1}}%
	{\XXint\textstyle\scriptstyle{#1}}%
	{\XXint\scriptstyle\scriptscriptstyle{#1}}%
	{\XXint\scriptscriptstyle\scriptscriptstyle{#1}}%
	\!\int}
\def\XXint#1#2#3{{\setbox0=\hbox{$#1{#2#3}{\int}$ }
		\vcenter{\hbox{$#2#3$ }}\kern-.6\wd0}}
\def\ddashint{\Xint=}
\def\dashint{\Xint-}

%\setlength{\oddsidemargin}{0.125in}
%\setlength{\evensidemargin}{0.125in}
%\setlength{\textwidth}{6.375in}
%\setlength{\textheight}{8.5in}
%\setlength{\extrarowheight}{0.2cm}
%\topskip 0in
%\topmargin 0.375in
%\footskip 0.25in
\renewcommand{\baselinestretch}{1.08}
%\numberwithin{equation}{section}
\renewcommand{\Re}{\operatorname{Re}}
\renewcommand{\Im}{\operatorname{Im}}
%\newcommand{\intprod}{\mathbin{\raisebox{\depth}{\scalebox{1}[-1]{$\lnot$}}}}
\DeclareMathSymbol{\intprod}{\mathbin}{MnSyC}{'270}
\newcommand{\LB}{\left[}
\newcommand{\RB}{\right]}
\newcommand{\LA}{\left\langle}
\newcommand{\RA}{\right\rangle}
\newcommand{\m}{{\rm modulo}\:}
\newcommand{\im}{{\rm im}\:}
\newcommand{\GI}{{\mathbb Z}[i]}
\newcommand{\mo}{{\rm mod}\:}
\newcommand{\F}{{\mathbb F}}
\newcommand{\Z}{{\mathbb Z}}
\newcommand{\N}{{\mathbb N}}
\newcommand{\C}{{\mathbb C}}
\renewcommand{\P}{{\mathbb P}}
\newcommand{\Q}{{\mathbb Q}}
\newcommand{\R}{{\mathbb R}}
\newcommand{\G}{{\mathbb G}}
\newcommand{\A}{{\mathbb A}}
\newcommand{\bbH}{{\mathbb H}}
\newcommand{\bi}{{\mathbf i}}
\newcommand{\bj}{{\mathbf j}}
\newcommand{\bk}{{\mathbf k}}
\newcommand{\rB}{{\mathrm{B} }}
\newcommand{\rH}{{\mathrm{H} }}
\newcommand{\SU}{{\mathrm{SU} }}
\newcommand{\Sp}{{\mathrm{Sp} }}
\newcommand{\U}{{\mathrm U}}
\newcommand{\SO}{{\mathrm{SO} }}
\renewcommand{\rB}{{\mathrm B}}
\renewcommand{\G}{{\mathrm G}}
\newcommand{\rK}{{\mathrm K}}
\newcommand{\rHol}{{\mathrm Hol }}
\newcommand{\Spin}{{\mathrm{Spin} }}
\newcommand{\supp}{{\mathrm{supp} \, }}
%\newcommand{\cO}{\mathcal{O}}
\renewcommand{\div}{{\mathrm{div} }}
\newcommand{\bbP}{{\mathbb P}}
\newcommand{\del}{{\partial}}
\newcommand{\delbar}{{\overline{\partial}}}
\newcommand{\gm}{{\mathfrak m}}
\newcommand{\gn}{{\mathfrak n}}
\newcommand{\gothg}{{\mathfrak g}}
\newcommand{\gothso}{\mathfrak {so }}
\newcommand{\gothk}{{\mathfrak k}}
\newcommand{\gothh}{{\mathfrak h}}
\newcommand{\gothu}{{\mathfrak u}}
\newcommand{\gothb}{{\mathfrak b}}
\newcommand{\gq}{{\mathfrak q}}
\newcommand{\gothC}{{\mathfrak C}}
\newcommand{\gothG}{{\mathfrak G}}
\newcommand{\gothB}{{\mathfrak B}}
\newcommand{\scrA}{{\mathscr A}}
\newcommand{\scrD}{{\mathscr D}}
\newcommand{\scrF}{{\mathscr F}}
\newcommand{\scrG}{{\mathscr G}}
\newcommand{\scrI}{{\mathscr I}}
\newcommand{\scrJ}{{\mathscr J}}
\newcommand{\scrL}{{\mathscr L}}
\newcommand{\scrO}{{\mathscr O}}
\newcommand{\scrM}{{\mathscr M}}
\newcommand{\T}{{\mathcal T}}
\newcommand{\Grass}{{\mbox{Grass}}}
\newcommand{\Spec}{{\mbox{Spec}}}
\newcommand{\Frob}{{\mbox{Frob}}}
\newcommand{\dist}{{\mbox{dist}}}
\newcommand{\Pic}{{\text{Pic}}}
\newcommand{\Div}{{\text{Div}}}
\newcommand{\End}{{\text{End}}}
\newcommand{\dbar}{{\bar{\partial}}}
\newcommand{\pd}{{\partial}}
\newcommand{\doublesup}[1][0pt]{%
	\mathrel{\raisebox{#1}{$\sup$}}%
}
\newcommand{\inj}{{\mathrm{inj} }}
\newtheorem{thm}{Theorem}[section]
\newtheorem{lemma}[thm]{Lemma}
\newtheorem*{lemma*}{Lemma}
\newtheorem{prop}[thm]{Proposition}
\newtheorem{ass}[thm]{Assumption}
\newtheorem{cor}[thm]{Corollary}
\newtheorem{conj}[thm]{Conjecture}
\newtheorem*{conj*}{Conjecture}
\newtheorem{defnlemma}[thm]{Definition/Lemma}
%\newtheorem{defn}{Definition}
%\newtheorem{example}{Example}
%\newtheorem{rmk}{Remark}
%\newenvironment{defn}[1][Definition.]{\begin{trivlist}
%\item[\hskip \labelsep {\bfseries #1}]}{\end{trivlist}}
%\newenvironment{example}[1][Example.]{\begin{trivlist}
%\item[\hskip \labelsep {\bfseries #1}]}{\end{trivlist}}
%\newenvironment{rmk}[1][Remark.]{\begin{trivlist}
%\item[\hskip \labelsep {\bfseries #1}]}{\end{trivlist}}

   \newtheoremstyle{others}% name
     {3pt}%      Space above
     {2pt}%      Space below
     {}%         Body font
     {}%         Indent amount (empty = no indent, \parindent = para indent)
     {\bf}% Thm head font
     {.}%        Punctuation after thm head
     {.5em}%     Space after thm head: " " = normal interword space;
           %       \newline = linebreak
     {}%         Thm head spec (can be left empty, meaning `normal')

\theoremstyle{others}
\newtheorem{rmk}[thm]{Remark}
\newtheorem*{rmk*}{Remark}
\newtheorem{defn}[thm]{Definition}
\newtheorem{example}[thm]{Example}
\newtheorem{exercise}[thm]{Exercise}


\numberwithin{equation}{section}

\setcounter{tocdepth}{1}

% ==Chong's customized commands============================
%Math operators
 \DeclareMathOperator{\tr}{tr}
 \DeclareMathOperator{\vol}{vol}
% Math Notations ---------------------------------------------------------
 \newcommand{\dif}[1]{\frac{d}{d{#1}}}
 \newcommand{\Real}{\mathbb{R}}
 \newcommand{\Integer}{\mathbb{Z}}
 \newcommand{\Complex}{\mathbb{C}}
 \newcommand{\Abs}[1]{\left\vert#1\right\vert}
 \newcommand{\abs}[1]{\vert#1\vert}
 \newcommand{\norm}[1]{\Vert#1\Vert}
 \newcommand{\Norm}[1]{\left\Vert#1\right\Vert}
 \def\<{\left\langle} \def\>{\right\rangle}
 \def\({\left(} \def\){\right)}
 % Greek -------------------------------------------------------------------
 \newcommand{\al}{\alpha}
 \newcommand{\be}{\beta}
 \newcommand{\ld}{\lambda}
 \newcommand{\Ld}{\Lambda}
 \newcommand{\de}{\delta}
 \newcommand{\De}{\Delta}
 \newcommand{\eps}{\varepsilon}
 \newcommand{\Si}{\Sigma}
 \newcommand{\si}{\sigma}
 \newcommand{\om}{\omega}
 \newcommand{\Om}{\Omega}
 \newcommand{\ga}{\gamma}
 \newcommand{\Ga}{\Gamma}
 \newcommand{\ka}{\kappa}
 \renewcommand{\th}{\theta}
 \newcommand{\Th}{\Theta}

 \renewcommand{\t}{\tilde}
 \renewcommand{\b}{\bar}
 \renewcommand{\o}{\overline}
 \newcommand{\p}{\partial}
 \newcommand{\n}{\nabla}
 \newcommand{\bp}{\b{\p}}
 \newcommand{\dd}[1]{\frac{d}{d{#1}}}
 \newcommand{\pp}[1]{\frac{\p}{\p{#1}}}

\newcommand{\E}{\mathcal{E}}



\begin{document}



\title[Harmonic map flow with strict type-II blowup]{Strict type-II blowup in harmonic map flow}
\author{Alex Waldron}
\address{University of Wisconsin, Madison}
\email{waldron@math.wisc.edu}


\begin{abstract} A finite-time singularity of 2D harmonic map flow will be called ``strictly type-II'' if the outer energy scale satisfies
$$\lambda(t) = O (T-t)^{\frac{1 + \alpha}{2} }.$$
We prove that the body map at a strict type-II blowup is H\"older continuous. This is relevant to a conjecture of Topping. %for any $\beta < \alpha.$ %H\"older continuous. % and necks do not form.
\end{abstract}

\maketitle



\thispagestyle{empty}


%\tableofcontents





\section{Introduction}

Let $M$ and $N$ be compact Riemannian manifolds. For differentiable maps $u: M \to N,$ we may define the Dirichlet energy:
$$\frac12 \int_M |d u|^2 \, dV.$$
Its downward gradient flow is given by
\begin{equation}\label{hm}
\frac{\pd u}{\pd t} = \mathcal{T}(u).
\end{equation}
Here, $\T(u)$ is the tension field of $u,$ a generalization of the Laplace-Beltrami operator to maps between manifolds.
%The flow (\ref{hm}) is known as \emph{harmonic map flow}, and has been studied since 1964 \cite{eellssampson} almost without interruption.
The evolution equation (\ref{hm}), known as \emph{harmonic map flow}, was introduced in 1964 by Eells and Sampson \cite{eellssampson} and has been studied since then almost without interruption.
%The evolution equation (\ref{hm}), known as \emph{harmonic map flow}, was introduced in 1964 by Eells and Sampson \cite{eellssampson} and has been studied since then almost without interruption.
%The flow (\ref{hm}), introduced in 1964 by Eells and Sampson \cite{eellssampson}, is known as \emph{harmonic map flow}. %It has been studied since 1964 almost without interruption.

%The flow (\ref{hm}), known as \emph{harmonic map flow}, has been studied almost without interruption since its introduction by Eells and Sampson \cite{eellssampson} in 1964.
%, who used it to construct a harmonic map in any homotopy class, assuming that $N$ has nonpositive sectional curvature.
%For $\dim(M) = 2,$ this functional is conformally invariant, and this is the case of greatest interest.

When $M$ has dimension two, the Dirichlet functional is conformally invariant; we shall be concerned exclusively with this case. Struwe \cite{struwehmsurfaces} constructed a global weak solution $u(t)$ %$u : M \times \LB 0 , \infty \right) \to N$
of (\ref{hm}) starting from any initial map of Sobolev class $W^{1,2}(M,N).$ The solution is smooth away from finitely many singular times, when isolated singularities (``bubbles'') may form. For a singular time $T < \infty,$ the limit
$$u(T) = \lim_{t \nearrow T} u(t)$$
exists weakly in $W^{1,2}$ and smoothly away from the bubbling set, %at a time $T,$ 
%{\it A priori}, $u(T)$ belongs only to $W^{1,2}(M,N).$
%The regularity of the body map near the bubble points is a central question in the field.
and is referred to as the \emph{body map}. %The weak solution is obtained by discarding the bubbles and restarting the flow with initial data $u(T).$

%Since $W^{1,2}$ is not a subspace of $C^0,$
Note that Struwe's construction leaves open the possibility that the body map will be discontinuous.
Topping \cite{toppingwindingbehavior} demonstrated that for certain target manifolds and initial data, $u(T)$ can indeed have an essential singularity at a bubble point.
At the same time, he conjectured that for well-behaved (specifically, real-analytic) metrics on the target,
the body map always extends continuously across the bubbling set. Topping's conjecture is the {\it sine qua non} for future geometric applications of harmonic map flow.
%If Topping's conjecture is true, 
%the bubbling phenomenon should not prevent 
%then we can expect
%Continuity of the body map is the {\it sine qua non} for any further geometric application of harmonic map flow. % outside the setting of a nonpositively curved target manifold \cite{}.
%beyond the case of a nonpositively curved target manifold \cite{}.
%for studying maps to general (real-analytic) target manifolds, in addition to nonpositively curved ones \cite{}.
%(\ref{hm}) will remain useful for studying maps to (real-analytic) target manifolds
%beyond the non-positively curved case \cite{}. %in addition to non-positively curved ones \cite{}. %, as well as non-positively curved ones \cite{eellssampson}. %probe the space of all continuous maps $M \to N.$ %even for positively curved $N.$ %as it is for negatively curved target manifolds \cite{}.

In previous joint work with C. Song \cite{songwaldron}, we established H\"older continuity of the body map when $N$ is compact K\"ahler with nonnegative holomorphic bisectional curvature and the energy of the initial map is near the holomorphic energy. %``half'' the initial energy of the initial map is small.
%An important point in the argument was to establish a bound on the outer energy scale of the form
The argument relied partly on establishing a bound on the outer energy scale of the form
\begin{equation}\label{stricttypeII}
\lambda(t) = O(T - t)^{\frac{1 + \alpha}{2}},
\end{equation}
where $0 < \alpha \leq 1.$
%which always holds for harmonic map flow in dimension two (see {\it e.g.} \cite{}). 
We refer to a singularity satisfying (\ref{stricttypeII})
as a \emph{strict type-II blowup}. Note that (\ref{stricttypeII}) is a refinement of the ordinary type-II blowup estimate for 2D harmonic map flow: %, which automatically holds for 2D harmonic map flow:
\begin{equation}\label{typeII}
\lambda(t) = o(T - t)^{\frac12}.
\end{equation}
%which automatically holds in dimension two---
For a proof of (\ref{typeII}), see %{\it e.g.}
\cite[Theorem 1.6v]{toppingwindingbehavior}. % for a proof that (\ref{typeII}) always holds for 2D harmonic map flow.



The strict type-II bound (with $\alpha = 1$) is most familiar from the rotationally symmetric setting. Angenent, Hulshof, and Matano \cite{angenenthulshofmatano} proved that the finite-time blowup first constructed by Chang, Ding, and Ye \cite{changdingye} occurs with rate $\lambda(t) = o(T - t).$ %Continuity of the body map is well established there:
Raphael and Schweyer \cite{raphaelschweyer} determined a large set of rotationally symmetric initial data that blows up under (\ref{hm}) with the precise rate
\begin{equation}\label{raphaelschweyerblowuprate}
\lambda(t) \sim \kappa \frac{T - t}{(\ln T - t )^2}.
\end{equation}
They also proved in this context that the body map is $W^{2,2},$ hence $C^\beta$ for each $\beta < 1.$\footnote{This also follows from the main theorem of \cite{songwaldron}.} Davila, Del Pino, and Wei \cite{daviladelpinowei} produced a larger set of examples with blowup rate (\ref{raphaelschweyerblowuprate}), whose body maps are continuous by construction. %On the other hand, (\ref{stricttypeII}) can fail in certain cases; most notably, Topping's example of discontinuous blowup is not strictly type-II. (See \cite{}, Theorem bla e.) %the elementary type-II rate (\ref{typeII}) necessarily holds, but


In another direction, Topping \cite{toppingholder} proved continuity of the body map if the Dirichlet energy is H\"older continuous as a function of time. One step in the proof was to establish that the strict type-II bound (\ref{stricttypeII}) follows from this assumption (see \cite[Lemma 2.2]{toppingholder}). Hence, the strict type-II bound with arbitrary exponent has appeared in previous work, although a much stronger assumption was required to obtain continuity of the body map.

%The bound (\ref{}) also figures in Topping's proof that H\"older continuity of the Dirichlet energy implies continuity of the body map, although this last assumption is stronger than 

Accordingly, the known examples of strict type-II blowup all have continuous body maps. On the other hand, the only known example with discontinuous body map, due to Topping, fails to be strictly type-II---see \cite[Theorem 1.14e]{toppingwindingbehavior}. Our main theorem confirms the implication, as follows. %that (\ref{stricttypeII}) implies continuity of the body map.

\begin{comment}

Whereas any solution of 2D harmonic map flow must obey the ordinary type-II blowup estimate $\lambda(t) = o(T - t)^{\frac12},$
%\begin{equation*}%\label{typeII}
%\lambda(t) = o(T - t)^{\frac12},
%\end{equation*}
the sharper bound (\ref{stricttypeII}) can fail in certain cases. Indeed, Topping's example of discontinuous blowup is not strictly type-II (see \cite{toppingwindingbehavior}, Theorems 1.6v and 1.14e). On the other hand, according to the above discussion, all known examples of strict type-II blowup have continuous body map.
%Based on this observation and the rotationally symmetric evidence,
%However, it is reasonable to hope that the body map can be discontinuous only if the strict type-II condition fails.
%This is the content of our main result.
The implication is confirmed by the following result. %that the body map can be discontinuous only if the strict type-II condition fails.
%The main result of this paper is that the strict type-II condition is sufficient by itself for continuity of the body map.
%From this point of view, the main result of this paper is not altogether surprising:
%This paper proves that the strict type-II condition is sufficient for Topping's conjecture.  %the strict type-II rate (\ref{stricttypeII}) is sufficient in itself for continuity of the body map.

\end{comment}


\begin{thm}\label{thm:intro} For any strict type-II blowup of harmonic map flow in dimension two, with $0 < \alpha \leq 1$ in (\ref{stricttypeII}),
the body map is $C^{ \frac{\alpha}{3} }.$
\end{thm}

%\begin{comment}

\noindent Our main technical result is Theorem \ref{thm:mainthm} below, %, on p. \pageref{thm:mainthm} below. 
and Corollary \ref{cor:detailedthm} gives the formal statement of Theorem \ref{thm:intro}. %, on p. \pageref{cor:detailedthm}.
%We also establish the ``no-neck'' property under a slightly stronger assumption on the blowup rate, which is realistic based on \cite{ }. (See Theorem \ref{}).

%\end{comment}

%The detailed statement and proof of this result appears in Section \ref{sec:mainthm}, beginning on p. \pageref{sec:mainthm} below.
The proof %of our main result %of H\"older continuity
depends on obtaining decay estimates for the differential of $u$ in the ``neck region,'' {\it i.e.}, the area near the singularity but outside the (vanishing) energy scale.
The required estimate on the angular component of $du$ is already known \cite[Lemma 5.4]{songwaldron}, % (using similar calculations to those of \cite{linwangenergyidentity,qingradial})
so it remains only to estimate the radial component.
We obtain a bound on the difference between the radial and angular components from a well-known identity (\ref{stressidentity}), giving an integral bound on the radial component under the flow. %, which gives an integral bound under the flow. %(measured by the stress-energy tensor).
 %Proposition \ref{gestimate}, %an evolution equation for the radial component, (\ref{radialevol}),
Using a specialized parabolic estimate (Proposition \ref{gestimate}) and a bootstrap argument, we are able to promote this integral bound to pointwise decay throughout the neck region. %This yields the required decay estimates on the pointwise norm of $du$ farther out in the neck. %which directly imply H\"older continuity of the body map.




\subsection*{Acknowledgement} The author is partially supported by NSF DMS-2004661.



\section{Technical results}

For an introduction to harmonic map flow, we refer the reader to \cite[\S 2]{songwaldron} or to the textbook of Lin and Wang \cite{linwangbook}.

Let $(M,g)$ and $(N,h)$ be Riemannian manifolds.
Given any smooth map $u : M \to N,$ we denote the pullback of $h$ to $u^*TN$ by $\LA \cdot, \cdot \RA,$ which we combine with $g$ on tensors.
The differential $du$ is a section of $T^*M \otimes u^*TN,$ with norm squared
$$|du|^2 = g^{ij} \LA \p_i u, \p_j u \RA.$$
%We write $\nabla$ for either the Levi-Civita connection on $M,$ the pullback to $u^*TN$ of the Levi-Civita connection on $N,$ or their combination on tensors.
The tension field is given by
$$\T(u) = \tr_g \n du.$$
Here, $\nabla$ is the Levi-Civita connection on $T^*M$ coupled with the pullback to $u^*TN$ of the Levi-Civita connection on $N.$

Suppose that $\dim(M) = 2.$ %Choose conformal coordinates $X^i$ %$(r,\theta)$
We write
\begin{equation}\label{conformalmetric}
g=\xi^2(dr^2+r^2d\th^2)
\end{equation}
for the given metric in conformal coordinates, where $\xi$ is a smooth function. By adjusting the conformal chart to second order, we may assume $d \xi (0) = 0.$ We may further assume
\begin{equation}\label{smallness1}
\left| \xi - 1\right| + r|d \xi| + r^2 |\nabla d \xi| \leq \xi_0 r^2
\end{equation}
for a constant $0 \leq \xi_0 \leq \frac12,$ after rescaling $g$ by a constant. This implies
\begin{equation}\label{rversusdist}
|r- \dist_g(x_0, \cdot)| \leq C \xi_0 r^2.
\end{equation}
In view of these bounds, the difference between conformal and geodesic coordinates will be of no significance for our results; we use conformal coordinates only for convenience.

Let
$$S_{ij} = \LA \p_i u , \p_j u \RA - \frac12 g_{ij} |du|^2$$
be the \emph{stress-energy tensor} of $u.$ This is a symmetric 2-tensor on $M,$ which satisfies
\begin{equation}\label{stressdiv0}
\nabla^i S_{ij} = \LA \T(u), \p_j u \RA.
\end{equation}
 %Here, $\nabla$ is the connection on $TM \otimes u^*TN$ induced by the Levi-Civita connection on $M$ and the pullback of the Levi-Civita connection on $N.$
For a derivation of (\ref{stressdiv0}), see \cite[\S 2.2]{songwaldron}.

%Contracting with a vector field $Y^j,$ we have
%\begin{equation}\label{stressdiv}
%\nabla^i (Y^j S_{ij}) = \LA \T(u), Y^j \p_j u \RA + \frac12 \left( \nabla^i Y^j + \n^j Y^i \right) S_{ij}.
%\end{equation}
The radial vector field $X^i$ in conformal coordinates is conformally Killing, {\it i.e.}
\begin{equation*}
\nabla^i X^j + \n^j X^i = \lambda g^{ij}
\end{equation*}
for a scalar function $\lambda.$ Contracting with $X^i$ in (\ref{stressdiv0}), we have
\begin{equation}\label{stressdiv}
\begin{split}
\nabla^i (X^j S_{ij}) & = \LA \T(u), X^j \p_j u \RA + \frac12 \left( \nabla^i X^j + \n^j X^i \right) S_{ij} \\
& =  \LA \T(u), X^j \p_j u \RA,
\end{split}
\end{equation}
since $g^{ij}S_{ij} = 0$ in dimension two.
Integrating (\ref{stressdiv}) over a disk $D_r$ in the conformal chart, and applying the divergence theorem, we obtain
\begin{equation}\label{stressidentity}
\int_{S^1_r} X^i X^j S_{ij} \, d\theta = \int_{D_r} \LA \T(u), X^j \p_j u \RA \, dV. %+ O(r^2).
\end{equation}
%Notice that in polar coordinates, we have
%$$X^i X^j S_{ij} = r^2|u_r|^2 - |u_theta|^2.$$
%Applying H\"older's inequality, this yields
%\begin{equation}\label{stressholder}
%\begin{split}
%\left| \frac12 \int_{S^1_r} \left( r^2 |u_r|^2 - |u_\theta|^2 \right) d\theta \right| & = \left| \int_{S^1_r} X^i X^j S_{ij} \, d\theta \right| \\
%& \leq \left( \int_{B_r} |\T(u)|^2 \, dV \int_{B_r} {r'}^2 |du|^2 \, dV_{r'} \right)^{\frac12} + O(r^2).
%\end{split}
%\end{equation}
This identity is well known from the theory of approximate harmonic maps, and will be used crucially below.












%\section{Evolution of $du$}\label{ss:evolutionofangular}



Next, we need the following parabolic estimates. For $0 \leq \nu \leq 1$ and $\mu > 1,$ let
$$\square_\nu = \p_t - \left( \p_r^2 + \frac{1}{r} \p_r - \frac{\nu^2}{r^2} \right),$$
$$\Delta_\mu =\p_r^2 + \frac{(\mu - 1)}{r} \p_r.$$
We have
\begin{equation}\label{squaretriangle}
\square_\nu \left( r^{\nu} y \right) = r^\nu \left( \p_t - \Delta_{2(\nu + 1)} \right) y.
\end{equation} 
Also notice that
\begin{equation}\label{squarerbeta}
\square_\nu \! r^\beta = \left( \nu^2 - \beta^2 \right) r^{\beta - 2}.
\end{equation}


\begin{lemma}\label{lemma:evolutions}
Let $u$ be a solution of (\ref{hm}) %on $U^1_\rho \times [0, T)$
with respect to a metric $g$ of the form (\ref{conformalmetric}-\ref{smallness1}). Suppose that for some $0 < \eta^2 < \eps_0,$ % R \leq R_1 \leq 1,$ and $0 \leq t_1 < T,$
 we have
\begin{equation}\label{evolutions:derivbounds}
r|du|+ r^2|\n du| + r^3|\n^2 du| + r^4|\n^3 du| \le \eta. %\quad \forall \,  x \in U_{R_1}^1.
\end{equation}
Then the angular energy
\begin{equation*}%\label{angularenergy}
f = f(u;r, t) := \sqrt{ \int_{ S_r^1 } |u_\th (r,\theta,t)|^2 d\th + \int_{ S_r^1 } |\nabla_\theta u_\th (r,\theta,t)|^2 d\th }
\end{equation*}
satisfies a differential inequality
\begin{equation}\label{evolutions:angular}
\square_\nu \! f \le C \xi_0 \eta,
\end{equation}
where $\nu = \sqrt{1 - C \eta}.$
%for $t = t_1$ and $R_1 \leq r \leq 1.$
The radial energy
\begin{equation}\label{radialenergy}
g = g(u;r, t) := \sqrt{ \int_{ S_r^1 } r^2 |u_r (r,\theta,t)|^2 d\th }
\end{equation}
satisfies
\begin{equation}\label{evolutions:radial}
\square_\nu \! \left(\frac{g}{r} \right) \le \frac{6f}{r^3} + \frac{ C \xi_0 \eta}{r}.
\end{equation}
Here, $\varepsilon_0$ depends on the geometry of $N,$ and $\xi_0$ is the constant of (\ref{smallness1}).
%where $C_N$ is a constant depending on the geometry of $N$.
\end{lemma}
\begin{proof} The proof is an elementary extension of prior calculations; see Appendix \ref{app:evolutions}.
\end{proof}



%To estimate the curvature components, we will need the following two lemmas.


\begin{prop}\label{festimate} Let $-\nu \leq \beta_i \leq \nu \leq 1,$ for $i = 0,1.$ %and put $\beta_0 = \frac12 \left( 3\beta - \nu \right).$
Suppose that $f(r,t)$ is continuous on $\LB \rho, 1 \RB \times \LB \tau, T \right)$ and satisfies
\begin{equation}\label{festimate:equation}
\Box_\nu \! f \leq A,
\end{equation}
with
\begin{equation}\label{festimate:assumptions}
\begin{split}
| f(r,\tau) | & \leq A \left( \left( \frac{\rho}{r} \right)^{\beta_0} + r^{\beta_1} \right), \\
f(\rho, t)  \leq & A, \qquad f(1, t)  \leq A.
\end{split}
\end{equation}
Then, given $0 < \kappa \leq 1/2,$ %and $\frac{\rho}{\kappa} \leq \rho_1 \leq \kappa,$
for
$$\frac{\rho}{\kappa} \leq r \leq \kappa \quad \text{and} \quad \tau + \frac{r^2}{\kappa^2} \leq t < T,$$
we have
\begin{equation}\label{festimate:conclusion}
| f(r,t) | \leq C_{\ref{festimate}} A \left( \kappa^{\nu + \beta_0} \left( \frac{\rho}{r} \right)^{\beta_0} + \kappa^{\nu - \beta_1} r^{\beta_1} \right).
\end{equation}
Here $C_{\ref{festimate}}$ depends on $\beta_0, \beta_1,$ and $\nu.$
\end{prop}
\begin{proof} Let $\mu = 2 \nu + 2.$ Using the results of Appendix \ref{app:initialkernel} and (\ref{squaretriangle}-\ref{squarerbeta}), we can construct a supersolution for (\ref{festimate:equation}) of the form
$$\bar{v}(r,t) = r^\nu  v_0(r,t) + 2 A \left( \left( \frac{\rho}{r} \right)^\nu + r^\nu \right) - \frac{A}{4 - \nu^2} r^{2},$$
where
$$\left( \p_t - \Delta_{\mu} \right) v_0 = 0,$$
$$v_0(r,0) = r^{-\nu} f(r,0), \qquad v_0(\rho,t) = 0 = v_0(1,t).$$
Applying the comparison principle to (\ref{festimate:equation}), we have
\begin{equation}\label{festimate:maxprin}
g(r,t) \leq \bar{v}(r,t),
\end{equation}
so it suffices to check (\ref{festimate:conclusion}) for $\bar{v}.$

By Proposition \ref{initialestimate} and (\ref{festimate:assumptions}), we have
\begin{equation*}
v_0(r,t) \leq C A \left( \frac{\rho^{\beta_0} }{r^{\beta_0 + \nu} }w^{\beta_0 + \nu} \left( r,t - \tau\right) + r^{\beta_1 - \nu} w^{\nu - \beta_1 }(r,t - \tau) \right),
\end{equation*}
where $w^a(r,t)$ is defined by
$$w^a(r,t) = \left( \frac{r^2}{r^2 + t} \right)^{a/2}.$$
Overall, from (\ref{festimate:maxprin}), we obtain
\begin{equation}\label{festimate:flongest}
| f(r,t) | \leq CA \left( \left( \frac{\rho}{r} \right)^{\beta_0} w^{\beta_0 + \nu}(r,t - \tau) + r^{\beta_1} w^{\nu - \beta_1}(r,t - \tau) + \left( \frac{\rho}{r} \right)^\nu  + r^\nu + r^{2} \right).
\end{equation}
For $r \geq \rho/ \kappa,$ we have
$$\left( \frac{\rho}{r} \right)^\nu \leq \kappa^{\nu - \beta_0} \left( \frac{\rho}{r} \right)^{\beta_0}.$$
For $r \leq \kappa$ and $t - \tau \geq r^2 / \kappa^2,$ we have
$$ w^{\beta_0 + \nu}(r,t - \tau) \leq \kappa^{\beta_0 + \nu} , \qquad w^{\nu - \beta_1}(r,t - \tau) \leq \kappa^{\nu - \beta_1}, \qquad r^\nu \leq \kappa^{\nu - \beta_1} r^{\beta_1}.$$
Substituting into (\ref{festimate:flongest}), we obtain (\ref{festimate:conclusion}).
\end{proof}



\begin{prop}\label{gestimate} Let $- \nu \leq \gamma_i \leq \nu \leq 1$ and $\beta_i$ with $|1 \pm \beta_i| \neq \nu,$ for $i = 0,1.$ Given $0 < \kappa \leq 1/2,$ let $\rho \leq \rho_1 \leq \kappa.$

Suppose that $g(r,t)$ is continuous on $\LB \rho_1, 1 \RB \times \LB \tau, T \right)$ and satisfies
\begin{equation}\label{gestimate:equation}
\square_\nu \! \left( \frac{g}{r} \right) \leq \frac{A}{r^3} \left( \left(\frac{\rho}{r}\right)^{\beta_0} + r^{\beta_1} \right),
\end{equation}
with
\begin{equation}\label{gestimate:assumptions}
\begin{split}
|g(r,\tau)| & \leq A \left( \left( \frac{\rho}{r}\right)^{\gamma_0} + r^{\gamma_1} \right), \\
\int_{\tau}^{T} \!\! g(\rho_1, t)^2\, dt & \leq B^2, \qquad | g(1,t) | \leq A.
\end{split}
\end{equation}
Then, %and $\frac{\rho}{\kappa} \leq \rho_1 \leq \kappa,$
for
\begin{equation}\label{gestimate:rtassumptions}
2 \rho_1 \leq r \leq \kappa \quad \text{and} \quad \tau + \frac{r^2}{\kappa^2} \leq t < T,
\end{equation}
we have
\begin{equation}\label{gestimate:estimate}
| g(r,t) | \leq C_{\ref{gestimate}} \left( B \frac{\rho_1^{\nu - 1}}{r^\nu}+ A \left( \kappa^{\nu + \gamma_0 + 1}  \left( \frac{\rho}{r} \right)^{\gamma_0} + \kappa^{\nu - \gamma_1 + 1} r^{\gamma_1} + \left( \frac{\rho}{r} \right)^{\beta_0} + r^{\beta_1} \right)\right).
\end{equation}
Here, $C_{\ref{gestimate}}$ depends on $\gamma_0, \gamma_1, \beta_0, \beta_1,$ and $\nu.$
\end{prop}
\begin{proof} As in the previous proof, we construct a supersolution for (\ref{gestimate:equation}) of the form
$$\bar{v}(r,t) = r^\nu \left( v_1(r,t) + v_2(r,t) \right) + C A \left( \frac{\rho^{\beta_0}}{r^{\beta_0 + 1}} + r^{\beta_1 - 1} \right),$$
where
$$\left( \p_t - \Delta_\mu \right) v_i = 0, \quad i = 1,2,$$
with
$$v_1(r,0) = r^{-\nu - 1} g(r,0), \qquad v_1(\rho,t) = 0 = v_1(1,t),$$
and
$$v_2(\rho_1, t) = \rho_1^{-\nu - 1} g(\rho_1, t), \qquad v_2(r,0) = 0= v_2(1,t).$$
Applying the comparison principle to (\ref{gestimate:equation}), we have
\begin{equation}\label{gestimate:maxprin}
g(r,t) \leq r \bar{v}(r,t).
\end{equation}
By Proposition \ref{innerestimate} and (\ref{gestimate:assumptions}), we have
\begin{equation*}
v_1(r,t) \leq C A \left( \frac{\rho_1^{\gamma_0}}{r^{\gamma_0 + \nu + 1} }w^{\gamma_0 + \nu + 1} \left( r,t \right) + r^{\gamma_1 - \nu - 1} w^{\nu - \gamma_1 + 1}(r,t) \right).
\end{equation*}
By Proposition \ref{innerestimate} and (\ref{gestimate:assumptions}), since $\mu - 2 = 2\nu,$ we have
\begin{equation*}
v_2(r,t) \leq C B \rho_1^{-\nu - 1} \left( \frac{\rho_1^{2 \nu}}{r^{2 \nu + 1}} \right) \leq CB \frac{\rho_1^{\nu - 1}}{r^{2\nu + 1}}.
\end{equation*}
Overall, from (\ref{gestimate:maxprin}), we obtain
\begin{equation*}
| g(r,t) | \leq C\left( \frac{B \rho_1^{\nu - 1}}{r^\nu} + A \left( \left( \frac{\rho_1}{r} \right)^{\gamma_0} w^{\gamma_0 + \nu + 1}(r,t) +  r^{\gamma_1} w^{\nu - \gamma_1 + 1}(r,t) + \left( \frac{\rho}{r} \right)^{\beta_0} + r^{\beta_1} \right) \right).
\end{equation*}
As in the previous proof, under the assumptions (\ref{gestimate:rtassumptions}) on $r$ and $t,$ this implies (\ref{gestimate:estimate}).
\end{proof}




\section{Main Theorem}\label{sec:mainthm}


Fix $x_0 \in M$ and let $r_g(x) = \dist_g(x_0, x).$ We shall denote the geodesic ball of radius $R$ centered at $x_0$ by $B_R(x_0),$ or more typically by $B_R.$ An annulus will be denoted by $U_\rho^R = B_R \setminus \bar{B}_\rho.$
%\begin{equation*}
 %U_\rho^1 = \{ (r, \theta) \mid \rho \leq r \leq 1 \}.
%\end{equation*}
%denote an annulus in the plane. 

\begin{defn}[\cite{songwaldron}, Definition 4.2] Given a $W^{1,2}$ map $u: B_R(x_0) \to N,$ the \emph{outer energy scale} $\lambda_{\eps, R, x_0}(u)$ is the smallest nonnegative number $\lambda$ such that
\begin{equation}\label{energyscale}
\sup_{\lambda < \rho < R} \int_{U^{\rho}_{\rho/2}(x_0)} |du|^2 \, dV < \eps.
\end{equation}
Note that $\lambda = R$ satisfies (\ref{energyscale}) vacuously, so $0 \leq \lambda_{\eps, R, x_0}(u) \leq R$ by definition.
\end{defn}
We first establish the following ``baby case'' of the main theorem. %a lemma which should be well known, but whose sharp form does not appear in the literature (to our knowledge). It is a ``baby case'' of the main theorem.



\begin{lemma}\label{basecase} Given $E, \lambda> 0$ and $0 < \eps < \eps_0,$ there exists $\delta > 0$ as follows. Suppose that $u$ is a smooth solution of (\ref{hm}) on $B_R \times ( - R^2, T),$ with $0 < R < R_0$ and $T > 0,$ which satisfies
\begin{equation}\label{basecase:energyassumption}
\sup_{-R^2<t<T} \int_{B_R} |du(t)|^2 \, dV \leq E,
\end{equation}
\begin{equation}\label{basecase:deltaassumption}
\int_{-R^2}^T \!\! \int_{B_R} |\T|^2 \, dV dt < \delta^2,
\end{equation}
and
\begin{equation}\label{basecase:lambdaassumption}
\sup_{-R^2<t<T} \lambda_{\eps, R, x_0}(u(t)) \leq R \lambda.
\end{equation}
Then for $2 R \lambda \leq r_g(x) \leq R/2$ and each integer $k \geq 0,$ we have
\begin{equation}\label{basecase:est}
r_g(x)^{1 + k} |\nabla^{(k)} du(x,0)| \leq C_k \sqrt{\eps} \left( \frac{R \lambda}{r_g(x)} + \frac{r_g(x)}{R} \right).
\end{equation}
Here, $R_0 > 0$ depends on the geometry of $M,$ $\eps_0 > 0$ depends on the geometry of $N,$ and $\delta$ depends on $E,\eps,$ and $\lambda.$
\end{lemma}
\begin{proof} 
Assuming that $R_0$ is sufficiently small, we may rescale so that $R = 1$ and our metric $g$ takes the form (\ref{conformalmetric}-\ref{smallness1}), with $\xi_0 \leq \frac12.$ By (\ref{rversusdist}), it suffices to establish (\ref{basecase:est}) with the conformal coordinate $r$ in place of $r_g.$



In view of (\ref{basecase:lambdaassumption}), the standard $\eps$-regularity lemma (see {\it e.g.} Theorem 3.4 of \cite{songwaldron}) implies
\begin{equation}\label{basecase:crudestimate}
\sup_{\lambda \leq r \leq R} \left( r |du| + r^2 |\nabla du| + r^3 |\nabla^2 du | + r^4 |\nabla^3 du | \right) \leq C \sqrt{\eps}.
\end{equation}
Hence, by Lemma \ref{lemma:evolutions}, $f$ satisfies an evolution equation
$$\square_\nu f \leq C \sqrt{\eps},$$
where $\nu = \sqrt{1 - C \sqrt{\eps}}.$
%By taking $R_0$ smaller, we can assume that $\xi_0$ is small, so this becomes
%$$\square_\nu f \leq \sqrt{\eps}.$$
We can apply Proposition \ref{festimate}  with $\rho = \lambda,$ $A = C \sqrt{\eps},$ $\beta_0 = \beta_1 = 0,$ and $\tau = -\frac12.$ Since $w^a(r, s) \leq C r^a$ for $s \geq \frac14,$ (\ref{festimate:flongest}) gives
\begin{equation}\label{basecase:improvedfestimate}
\begin{split}
\sup_{-\frac14 \leq t \leq 0} | f(r,t) | & \leq C\sqrt{\eps} \left( r^{\nu} + r^{\beta_1} r^{\nu - \beta_1} + \left( \frac{\lambda}{r} \right)^\nu + r^\nu + r^{2} \right) \\
& \leq C \sqrt{\eps} \left( \left( \frac{\lambda}{r} \right)^\nu  + r^\nu \right).
\end{split}
\end{equation}
%Now, according to Theorem 3.4 of \cite{songwaldron}, we have a pointwise bound on the tension field
%\begin{equation}
%|\T(u(x,t)) | \leq \frac{C}{\lambda^2} \delta
%\end{equation}
%for $\lambda \leq |x| \leq 1$ and $-\frac12 \leq t \leq 0.$ 
Applying H\"older's inequality identity to (\ref{stressidentity}), we have
\begin{equation}\label{basecase:holder}
\begin{split}
\left| \int_{-\frac12}^{0} \left( g^2(r, t) - f^2(r, t) \right) \, dt \right| % & = \left|  \int_{t_0}^{t_1} \int_{S^1_r} \left( r^2 |u_r|^2 - |u_\theta|^2 \right) d\theta dt \right| \\
& = \left| \int_{t_0}^{t_1} \!\!\!\!\! \int_{S^1_{\rho_1}} X^i X^j S_{ij} (x,t) \, d\theta dt \right| \\
& \leq \left( \int_{t_0}^{t_1} \!\!\!\!\! \int_{D_{\rho_1}} |\T(u)|^2 \, dV dt  \right)^{\frac12} \left( \int_{t_0}^{t_1} \!\!\!\!\!  \int_{D_{\rho_1}} r'^2 \, |du|^2 \, dV_{r'} dt \right)^{\frac12} \\
& \leq C \delta \sqrt{E}.
\end{split}
\end{equation}
Since $|du|^2 \leq f^2 + g^2,$ for $\delta$ sufficiently small, (\ref{basecase:improvedfestimate}-\ref{basecase:holder}) imply
\begin{equation*}
\begin{split}
\int_{-\frac12}^0 \int_{r/2}^r |du|^2 \, dV dt & \leq C \eps \left( \left( \frac{\lambda}{r} \right)^\nu  + r^\nu \right)^2.
\end{split}
\end{equation*}
Applying $\eps$-regularity (Theorem 3.4 of \cite{songwaldron}) again, we obtain
\begin{equation}\label{basecase:improvedestimate}
\begin{split}
r^{1 + k} |\nabla^{(k)} du(x,0)| & \leq C_k \sqrt{\eps} \left( \left( \frac{\lambda}{r} \right)^\nu  + r^\nu \right).
\end{split}
\end{equation}
%which is almost the desired estimate (\ref{basecase:est}).
To obtain the same estimate with $\nu = 1,$ one can apply a supersolution argument as in Propositions \ref{festimate}-\ref{gestimate}, with (\ref{basecase:improvedestimate}) in place of (\ref{basecase:crudestimate}). Since the sharp result will not be used below, we omit the proof. %This yields (\ref{basecase:improvedfestimate}) with $\nu = 1.$
%Choosing $\delta$ smaller, one obtains the same estimate for $g.$
\end{proof}





\begin{thm}\label{thm:mainthm} %$0 < \lambda_1 \leq \lambda_0,$ and
Given $E > 0,$ $0 < \eps < \eps_0,$ and $0 < \alpha \leq 1,$ there exists $\delta_0 > 0$ as follows.  %and $\frac12 < \nu \leq \sqrt{1 - C \eps}.$ %and choose $\gamma$ with
%\begin{equation}\label{mainthm:gammaassumption}
%0 < \gamma < \min \LB \beta, \frac14 \left( 2 \nu - 1 \right) \RB.
%\end{equation}
Suppose that $u : B_R \times \left( -R^2, T \right) \to N$ is a smooth solution of (\ref{hm}), with $0 < R < R_0$ and $T > 0,$ which satisfies
\begin{equation}\label{mainthm:energyassumption}
\sup_{-R^2<t<T} \int_{B_R} |du(t)|^2 \, dV \leq E
\end{equation}
and
\begin{equation}\label{mainthm:deltaassumption}
\int_{-R^2}^T \!\! \int_{B_R} |\T|^2 \, dV dt < \delta_0^2.
\end{equation}
Suppose further that for some $0 < \lambda_1 \leq \frac12$ and $0 \leq t_1 < T,$ we have
\begin{equation}\label{mainthm:lambdaassumption}
\lambda_{\eps, R, x_0}(u(t)) \leq R \left( \lambda_1 + \left( \frac{t_1 - t}{R^2} \right)^{\frac{1 + \alpha}{2}} \right)
\end{equation}
for all $-R^2 < t \leq t_1.$ %Suppose further that (\ref{stricttypeII}) is satisfied, for some $t_1 < T,$ with $0 < \alpha \leq \frac12,$ and let $0 < \beta < \alpha.$
Then for $2 R \lambda_1 \leq r_g(x) \leq R / 2$ and each integer $k \geq 0,$ we have
\begin{equation}\label{mainthm:est}
r_g(x)^{1 + k} |\nabla^{(k)} du (x,t_1)| \leq C_{k,\alpha} \sqrt{\eps} \left( \frac{R\lambda_1}{r_g(x)} + \left( \frac{r_g(x)}{R} \right)^\alpha \right)^{\frac13 }.
\end{equation}
Here, $R_0 > 0$ depends on the geometry of $M,$ $\eps_0 > 0$ depends on the geometry of $N,$ and $\delta_0$ depends on $E,\eps,$ and $\alpha.$ %and $C_\alpha$ depends only on $\alpha.$ % and $\nu.$
\end{thm}
\begin{proof} %Let $\beta = 2 \nu - 1$ and $\gamma = \frac14 \beta.$
For convenience, we replace (\ref{mainthm:lambdaassumption}) by
\begin{equation}\label{mainthm:kappathing}
\lambda_{\eps, R, x_0}(u(t)) \leq R \left( \lambda_1 + \kappa^{a} \left( \frac{t_1 - t}{R^2} \right)^{\frac{1 + \alpha}{2}} \right),
\end{equation}
where $0 < \kappa \leq \frac12$ is a constant % (depending only on $\alpha$ and $\nu$) 
to be determined by the argument below (depending only on $\alpha$ and $\nu$), and $a$ is sufficiently large, for instance
\begin{equation}\label{mainthm:athing}
a = \frac{72( 1 + \alpha) }{\alpha}.
\end{equation}
The assumptions (\ref{mainthm:kappathing}) and (\ref{mainthm:lambdaassumption}) are equivalent after rescaling and redefining constants. %implies the estimate (\ref{mainthm:est}).

We will also prove the two estimates
\begin{equation}\label{mainthm:fhypoth}
 f(u;r,t_1) \leq C_0 \sqrt{\eps} \max \LB \frac{\lambda_1}{r}, r^{\alpha} \RB^{2 \nu - 1},
\end{equation}
\begin{equation}\label{mainthm:ghypoth}
g(u;r,t_1) \leq C_0 \sqrt{\eps} \max \LB \frac{\lambda_1}{r}, r^{\alpha} \RB^{\frac13},
\end{equation}
which clearly imply (\ref{mainthm:est}). Here, $\nu$ is any number with
\begin{equation}\label{mainthm:nucloseness}
\frac{26}{27} \leq \nu \leq \sqrt{1 - C \sqrt{\eps}} .
\end{equation}
It suffices to prove the theorem for $\lambda_1$ of the form
$$\lambda_1 = 2^{-n} \kappa^a,$$
so we may proceed by induction. For $\lambda_1 = \kappa^a,$
Lemma \ref{basecase} gives $\delta_0 = \delta > 0$ such that (\ref{mainthm:fhypoth}-\ref{mainthm:ghypoth}) hold, with $C_0 > 1$ universal. This establishes the base case. Note that we are free to assume $\kappa$ is arbitrarily small.

%Note that by the base case, we have the freedom to assume that $\lambda_0$ is as small as we wish (at the price of taking $\delta_0$ smaller). 

For the induction step, suppose that (\ref{mainthm:fhypoth}-\ref{mainthm:ghypoth}) hold for all $\lambda_1 \geq 2\bar{\lambda}_1,$ where $0 < \bar{\lambda}_1 \leq \kappa^a;$ {\it i.e.}, the conclusion of the theorem holds for all such $\lambda_1$ and solutions $u$ satisfying the hypotheses. We must establish the Theorem for $\lambda_1 = \bar{\lambda}_1.$

Let $u(t)$ be a solution satisfying the hypotheses, with $\lambda_1 = \bar{\lambda}_1.$ By rescaling, it suffices to assume $R = 1.$


In view of (\ref{mainthm:lambdaassumption}), %for $\delta_0$ sufficiently small,
the standard $\eps$-regularity lemma (see {\it e.g.} Theorem 3.4 of \cite{songwaldron}) implies
$$\sup_{\lambda(t) \leq r \leq R} \left( r |du| + r^2 |\nabla du| + r^3 |\nabla^2 du | + r^4 |\nabla^3 du |\right) \leq C \sqrt{\eps}.$$
Hence, by Lemma \ref{lemma:evolutions}, $f$ and $g$ satisfy evolution equations
\begin{equation*}
\begin{split}
\square_\nu f & \leq C \sqrt{\eps} \\
\square_\nu \left( \frac{g}{r} \right) & \leq \frac{6f}{r^3} + C \frac{\sqrt{\eps} }{r},
\end{split}
\end{equation*}
for $\lambda_1 + \kappa^{a} \left(t_1 - t \right)^{\frac{1 + \alpha}{2}} \leq r \leq 1.$ 
Let
\[
\begin{split}
\rho = 2 \bar{\lambda}_1, \qquad \zeta  = \rho^{\frac{1}{1 + \alpha}}.
\end{split}
\]
Since $\bar{\lambda}_1 \leq \kappa^a$ (by the base case), we have
\begin{equation}\label{mainthm:zetasmallness}
\frac{\rho}{\kappa} \leq \zeta \leq 2\kappa^{\frac{a}{1 + \alpha}}.
\end{equation}
Also notice that $\rho = \zeta^{1 + \alpha},$ so
\begin{equation}\label{mainthm:rhoandzeta}
\frac{\rho}{\zeta} = \zeta^\alpha.
\end{equation}
We now apply the induction hypotheses to $u$ with $R = 1/2,$ $\lambda_1 = \rho,$ and $R \rho = \bar{\lambda}_1.$ %Since $f$ and $g$ are scale-invariant, 
In view of (\ref{mainthm:rhoandzeta}), we clearly have (\ref{mainthm:fhypoth}-\ref{mainthm:ghypoth}) for $r \leq \zeta / 2.$ Applying the induction hypothesis again, with $R = 1$ and $\lambda_1 = \rho,$ we also obtain (\ref{mainthm:fhypoth}-\ref{mainthm:ghypoth}) for all $r \geq 2 \zeta.$ 
Hence, it remains only to establish (\ref{mainthm:fhypoth}-\ref{mainthm:ghypoth}) for 
\begin{equation}\label{mainthm:rrange}
\frac12 \zeta \leq r \leq 2 \zeta.
\end{equation}
In other words, for $r$ as in (\ref{mainthm:rrange}), we must show
\begin{equation}\label{mainthm:frrangetoshow}
f(r,t_1) \leq \frac{C_0}{2} \sqrt{\eps} \zeta^{\alpha (2 \nu -1)}
\end{equation}
and
\begin{equation}\label{mainthm:grrangetoshow}
g(r,t_1) \leq \frac{C_0}{2} \sqrt{\eps} \zeta^{\frac{\alpha}{3}}.
\end{equation}
Let
$$t_0 = t_1 - \frac{\zeta^2}{\kappa^2}.$$
From (\ref{mainthm:kappathing}), we have $\lambda(t) \leq \rho$ for all $t_0 \leq t \leq t_1.$ By the induction hypothesis, we have
\begin{equation}\label{mainthm:fappliedindhypoth}
f(r,t) \leq C \sqrt{\eps} \left( \frac{\rho}{r} + r^{\alpha} \right)^{2 \nu - 1}
\end{equation}
\begin{equation}\label{mainthm:gappliedindhypoth}
g(r,t) \leq C \sqrt{\eps} \left( \frac{\rho}{r} + r^{\alpha} \right)^{\frac13}
\end{equation}
for all $t_0 \leq t \leq t_1,$ where $C$ is a multiple of $C_0.$
Combining these, we also have
\begin{equation}\label{mainthm:energydecayhypoth}
r |du(x,t)| \leq C \sqrt{\eps} \left( \frac{\rho}{r} + r^{\alpha} \right)^{\frac13}
\end{equation}
for all $t_0 \leq t \leq t_1.$

To obtain the estimate (\ref{mainthm:frrangetoshow}) on $f,$ we apply Proposition \ref{festimate}. From (\ref{mainthm:fappliedindhypoth}-\ref{mainthm:gappliedindhypoth}), we obtain
\begin{equation*}
\begin{split}
\sup_{\zeta/2 \leq r \leq 2 \zeta} f(r,t_1) & \leq C C_{\ref{festimate}} \sqrt{\eps} \left( \kappa^{1 - \nu} \left( \frac{\rho}{\zeta} \right)^{2\nu - 1} + \kappa^{\nu - \alpha (2\nu - 1)} \zeta^{\alpha (2\nu - 1)} \right) \\
& \leq C \sqrt{\eps} \left(\kappa^{1 - \nu} + \kappa^{\nu - \alpha (2\nu - 1)} \right)\zeta^{\alpha (2\nu - 1)},
\end{split}
\end{equation*}
where we have used (\ref{mainthm:zetasmallness}-\ref{mainthm:rhoandzeta}).
%Hence, assuming that $\kappa$ is sufficiently small, (\ref{mainthm:frrangetoshow}) is satisfied at $t = t_1.$
Assuming that $\kappa$ is small enough that
$$C \sqrt{\eps} \left(\kappa^{1 - \nu} + \kappa^{\nu - \alpha (2\nu - 1)} \right) \leq \frac{C_0}{2},$$
this establishes the desired estimate (\ref{mainthm:frrangetoshow}) on $f.$\footnote{The decay estimate on $f$ can also be obtained directly from \cite{songwaldron}, Lemma 5.4. We have re-proven it here (by a different method) for the sake of exposition.}

Next, to obtain the estimate on $g,$ let
$$\rho_1 = \zeta^{1 + \frac{9}{16} \alpha}.$$
By the induction hypothesis, we have
\begin{equation*}
\sup_{t_0 \leq t \leq t_1} f(\rho_1, t) \leq C \sqrt{\eps} \left(\frac{\rho }{\rho_1 } \right)^{2\nu - 1} \leq C \sqrt{\eps}  \zeta^{\frac{7\cdot 25}{16 \cdot 27}\alpha } \leq C \sqrt{\eps} \zeta^{\frac{3}{8} \alpha}.
\end{equation*}
Using (\ref{mainthm:nucloseness}), we obtain
\begin{equation}\label{mainthm:ftimeest}
\int_{t_0}^{t_1} f(\rho_1, t)^2 \, dt %\leq \frac{\zeta^2}{\kappa^2} C \eps \zeta^{\frac{7 \cdot 25}{8 \cdot 27} \alpha } 
\leq \frac{C \eps }{\kappa^2} \zeta^{2 + \frac{3}{4}\alpha }.
\end{equation}
We now integrate (\ref{stressidentity}) in time and apply H\"older's inequality:
\begin{equation}\label{mainthm:gfestimate1}
\begin{split}
\left| \int_{t_0}^{t_1} \left( g^2(\rho_1, t) - f^2(\rho_1, t) \right) \, dt \right| % & = \left|  \int_{t_0}^{t_1} \int_{S^1_r} \left( r^2 |u_r|^2 - |u_\theta|^2 \right) d\theta dt \right| \\
& = \left| \int_{t_0}^{t_1} \!\!\!\!\! \int_{S^1_{\rho_1}} X^i X^j S_{ij} (x,t) \, d\theta dt \right| \\
& \leq \left( \int_{t_0}^{t_1} \!\!\!\!\! \int_{D_{\rho_1}} |\T(u)|^2 \, dV dt  \right)^{\frac12} \left( \int_{t_0}^{t_1} \!\!\!\!\!  \int_{D_{\rho_1}} r'^2 \, |du|^2 \, dV_{r'} dt \right)^{\frac12} \\
& \leq C \delta_0 \left( (t_1 - t_0) \left( \rho^2 E + C \eps \rho_1^2 \left( \frac{\rho}{\rho_1} \right)^{\frac{2}{3} } \right) \right)^{\frac12}.
\end{split}
\end{equation}
Here we have used the assumptions (\ref{mainthm:energyassumption}-\ref{mainthm:deltaassumption}) and (\ref{mainthm:energydecayhypoth}).
We have
$$(t_1 - t_0) \rho^2 E = \frac{\zeta^{4 + 2 \alpha}}{\kappa^2} E$$
and 
$$(t_1 - t_0) \rho_1^2 \left( \frac{\rho}{\rho_1} \right)^{\frac{2}{3} \alpha} = \frac{\zeta^2}{\kappa^2} \zeta^{2 + \frac98 \alpha} \zeta^{\frac{7}{16} \frac23 \alpha} = \frac{ \zeta^{4 + \frac{17}{12} \alpha} }{\kappa^2}.$$
Hence, for $\delta_0$ sufficiently small (independently of $\bar{\lambda}_1$), (\ref{mainthm:gfestimate1}) reduces to
\begin{equation*}
\begin{split}
\left| \int_{t_0}^{t_1} \left( g^2(\rho_1, t) - f^2(\rho_1, t) \right) \, dt \right| % & = \left|  \int_{t_0}^{t_1} \int_{S^1_r} \left( r^2 |u_r|^2 - |u_\theta|^2 \right) d\theta dt \right| \\
& \leq \kappa \eps \zeta^{2 + \frac{17}{24} \alpha}.
\end{split}
\end{equation*}
Combining this with (\ref{mainthm:ftimeest}), we obtain
\begin{equation*}
\begin{split}
\int_{t_0}^{t_1} g(\rho_1, t)^2 \, dt & \leq \kappa \eps \zeta^{2 + \frac{17}{24} \alpha } + \frac{C \eps}{\kappa^2} \zeta^{2 + \frac{3}{4}\alpha } \\
& \leq C\eps  \left( \kappa + C \kappa^{\frac{1}{24}\frac{a \alpha}{1 + \alpha} - 2} \right) \zeta^{2 + \frac{17}{24} \alpha } \\
& \leq C \kappa \eps \zeta^{2 + \frac{17}{24}\alpha},
\end{split}
\end{equation*}
where we have used (\ref{mainthm:athing}) and (\ref{mainthm:zetasmallness}).
%By (\ref{mainthm:athing}), this reduces to
%\begin{equation}
%\begin{split}
%\int_{t_0}^{t_1} g(\rho_1, t)^2 \, dt & \leq C \eps \zeta^{2 + \frac{\alpha}{2}}.
%\end{split}
%\end{equation}

We may now apply Proposition \ref{gestimate}, to obtain
\begin{equation*}
\begin{split}
\sup_{\zeta/2 \leq r \leq 2 \zeta} g(r,t_1) & \leq C C_{\ref{gestimate}} \left( \begin{split}
& \frac{\sqrt{\kappa \eps} \zeta^{1 + \frac{17}{48} \alpha } }{\zeta^{1 + \frac{9}{16} \alpha} } \left( \frac{\zeta^{1 + \frac{9}{16}\alpha} }{\zeta} \right)^\nu + \sqrt{\eps} \left( \kappa^{\nu + \frac13 + 1} \left( \frac{\zeta^{1 + \alpha}}{\zeta}\right)^{\frac13} + \kappa^{\nu - \frac{\alpha}{3} + 1} \zeta^{\frac{\alpha}{3} } \right) \\
& \quad + \sqrt{\eps} \left( \frac{\zeta^{1 + \alpha} }{\zeta} + \zeta^{\alpha} \right)^{2 \nu - 1}
\end{split} \right) \\
& \leq C \left( \sqrt{\kappa \eps} \zeta^{\frac{\alpha}{48} \left( 17 + 27(\nu - 1) \right)} + \sqrt{\eps} \kappa \zeta^{ \frac{\alpha}{3} } + \sqrt{\eps} \zeta^{\frac{25}{27}\alpha} \right) \\
& \leq C \sqrt{\eps} \left( \sqrt{\kappa} + \kappa + \kappa^{\alpha \left( \frac{25a }{27(1 + \alpha)} - \frac13 \right)} \right) \zeta^{ \frac{\alpha}{3} },
\end{split}
\end{equation*}
where we have used (\ref{mainthm:nucloseness}) and (\ref{mainthm:zetasmallness}).
For $\kappa$ sufficiently small, this implies the desired bound (\ref{mainthm:grrangetoshow}), completing the induction.
\end{proof}


\begin{cor}[{\it Cf.} Theorem \ref{thm:intro}]\label{cor:detailedthm} Let $M$ be any Riemannian surface and suppose that $u$ is a classical solution of (\ref{hm}) on $M\times \LB 0, T \right)$
with bounded energy. Given $x_0 \in M$ and $0 < \eps < \eps_0,$ choose $0 < R < \min \LB R_0, \inj_{x_0}(M) \RB$ small enough that
\begin{equation}\label{lambdato0}
\qquad \lambda_{\eps, R, x_0} (u(t)) \to 0 \quad (t \nearrow T).
\end{equation}
If also
\begin{equation}\label{lambdaalphabound}
\lambda_{\eps, R, x_0} (u(t)) = O(T - t)^{\frac{1 + \alpha}{2}}, \qquad (t \nearrow T),
\end{equation}
where $0 < \alpha \leq 1,$ then $u(T)$ is $C^{\frac{\alpha}{3} }$ on $B_{R/2}(x_0).$
\end{cor}
\begin{proof} First note that since the energy is bounded, the stress-energy tensor is also bounded in $L^1.$ By Corollary 4.5 of \cite{songwaldron}, with $q = 1,$ (\ref{lambdato0}) is true for $R > 0$ sufficiently small.

Now, given any $\lambda_1 > 0,$ in view of (\ref{lambdaalphabound}), (\ref{mainthm:lambdaassumption}) will be satisfied for all $t_1$ sufficiently close to $T.$  %(shrinking $R$ further if necessary).
For $x \in U_{2\lambda_1}^R(x_0),$ by Theorem \ref{thm:mainthm}, we have
\begin{equation*}
r|du(x,T)| = \lim_{t_1 \nearrow T} r|du(x,t_1)| \leq C \sqrt{\eps} \left( \frac{R\lambda_1}{r} + \left( \frac{r}{R} \right)^\alpha \right)^{\frac13}.
\end{equation*}
Letting $\lambda_1 \searrow 0,$ we obtain
\begin{equation*}
r|du(x,T)| \leq C\sqrt{\eps} \left( \frac{r}{R} \right)^{ \frac{\alpha}{3}}
\end{equation*}
for all $x \in B_{R/2} \setminus \{x_0\}.$
Integrating radially, this gives $u(T) \in C^{\frac{\alpha}{3}}.$ % Taking $\eps$ smaller, we can make the same argument for $\nu < 1$ arbitrary. This completes the proof.
\end{proof}


\appendix





\section{Proof of Lemma \ref{lemma:evolutions}}\label{app:evolutions}



Let
\begin{equation*}%\label{angularenergy}
f_0(r) = \sqrt{ \int |u_\th (r,\theta,t)|^2 d\th}, \qquad f_1(r) = \sqrt{ \int |\nabla_\theta u_\th (r,\theta,t)|^2 d\th }.
\end{equation*}
In our previous paper \cite{songwaldron}, Lemma 5.1, we calculated the desired evolution of $f_0;$ we now apply a similar analysis to $f_1.$ These calculations go back to Lin-Wang \cite{linwangenergyidentity} and Parker \cite{parkerhmbubbletree}.



For convenience, we shall work below in cylindrical coordinates $(s=\ln r,\th).$ Letting
$$g_0=ds^2+d\th^2$$
denote the flat cylindrical metric, we have $g=\xi^{2} e^{2s} g_0.$
The differential of $u$ is given by
\begin{equation*}
du = u_s ds + u_\theta d\theta.
\end{equation*}
The tension field with respect to $g_0$ is given by
\[ \mathcal{T}_0(u)=\n_s u_s+\n_\th u_\th,\]
where $\n$ denotes the pullback connection on $u^*TN,$ as above.
The heat-flow equation (\ref{hm}) with respect to the metric $g$ becomes
\begin{equation}\label{e:hm-cylin}
u_t = \mathcal{T}(u) = \xi^{-2}e^{-2s}\mathcal{T}_0(u).
\end{equation}
%Define the angular energy of $u$ by
%\begin{equation}\label{angularenergy}
%f(u;r, t) := \sqrt{ \int_{ \{r\} \times S^1 } |u_\th (r,\theta,t)|^2 d\th }.
%\end{equation}
%For a subset $\Om'\subset (0,1]\times [0,T)$, let
%\[\Om=\{(x, t)\in B_1(p)\times [0,T)|(r,t)\in \Om'\}.\]
%First we deduce an evolution equation of the angular energy $f^2$, which is parallel to the corresponding ODE satisfied by a harmonic map (see for example Lemma 3.2 in~\cite{Parker1996}).
We start from the identity
\begin{equation}\label{ds2f12}
\frac12 \p_s^2f_1^2= \int_{S^1}\( |\n_s \nabla_\theta u_{\th}|^2+\< \n_s^2 \nabla_\theta u_{\th}, \nabla_\theta u_\th\>\).
\end{equation}
We have
\begin{equation}\label{T0ut}
\nabla_s u_s + \nabla_\theta u_\theta = \T_0(u) =  \xi^2 e^{2s} u_t.
\end{equation}
Applying $\n_\theta,$ we obtain
\begin{equation*}
\begin{split}
\nabla_\theta \nabla_s u_s + \nabla^2_\theta u_\theta & = 2 \xi \nabla_\theta \xi e^{2s} u_t + \xi^2 e^{2s} \nabla_\theta u_t \\
& =2 \xi \nabla_\theta \xi e^{2s} u_t + \xi^2 e^{2s} \nabla_t u_\theta,
\end{split}
\end{equation*}
and
\begin{equation}\label{dtheta2sus}
\begin{split}
\nabla_\theta^2 \nabla_s u_s & = - \nabla^3_\theta u_\theta + \left( 2 (\nabla_\theta \xi)^2 + 2 \xi \nabla^2_\theta \xi \right) e^{2s} u_t \\
& \quad + 2 \xi \nabla_\theta \xi e^{2s} \nabla_t u_\theta + \xi^2 e^{2s} \nabla_\theta \nabla_t u_\theta \\
& = - \nabla^3_\theta u_\theta + e^{2s} \nabla_t \nabla_\theta u_\theta + I,
\end{split}
\end{equation}
where
\begin{equation}\label{I}
I = (\xi^2 - 1) e^{2s} \nabla_\theta^2 u_t + e^{2s} R(u_\theta, u_t) u_\theta + \left( 2 (\nabla_\theta \xi)^2 + 2 \xi \nabla^2_\theta \xi \right) e^{2s} u_t + 2 \xi \nabla_\theta \xi e^{2s} \nabla_\theta u_t.
\end{equation}
We may also commute derivatives to obtain
\begin{equation}\label{ds2thetautheta}
\begin{split}
\nabla_s^2 \nabla_\theta u_\theta & = \nabla_s \left( \nabla_\theta \nabla_s u_\theta + R(u_s, u_\theta) u_\theta \right) \\
& = \nabla_\theta \nabla_s^2 u_\theta + \nabla_s \left( R(u_s, u_\theta) u_\theta \right) \\
& = \nabla_\theta \nabla_s ( \nabla_\theta u_s ) + \nabla_s \left( R(u_s, u_\theta) u_\theta \right) \\
& = \nabla_\theta^2 \nabla_s u_s + I \! I,
\end{split}
\end{equation}
where
\begin{equation}\label{II}
\begin{split}
I\! I & = \nabla_\theta \left( R(u_s, u_\theta) u_s \right) + \nabla_s \left( R(u_s, u_\theta) u_\theta \right) \\
& = \nabla R \left( u_\theta, u_s, u_\theta \right) u_s + R(\nabla_\theta u_s, u_\theta) u_s +R( u_s, \nabla_\theta u_\theta) u_s + R(u_s, u_\theta) \nabla_\theta u_s \\
& \quad + \nabla R \left( u_s, u_s, u_\theta \right) u_\theta + R(\nabla_s u_s, u_\theta) u_\theta + R( u_s, \nabla_s u_\theta) u_\theta + R( u_s, u_\theta) \nabla_s u_\theta.
\end{split}
\end{equation}
Inserting (\ref{dtheta2sus}) and (\ref{ds2thetautheta}) into (\ref{ds2f12}), integrating by parts, and rearranging, we obtain
\begin{equation}\label{firstevol}
\frac12 \left( e^{2s} \partial_t f_1^2 - \p_s^2f_1^2 \right) = - \int_{S^1}\( |\n_s \nabla_\theta u_{\th}|^2 + |\n_\theta^2 u_\theta|^2 \right) - \int_{S^1} \LA I + I\! I, \nabla_\theta u_\theta \RA.
\end{equation}
We need the following simple estimates. Since $|u_\th|\le \eta$ is small, we may assume that the image of the curve $u(s, \cdot, t):S^1\to N$ lies in a coordinate chart of $N$ where the Christoffel symbol $\Ga$ is bounded by $C_N$. Then, in local coordinates, we have
\[ |\p_\th u_\th| = |\n_\th u_\th - \Ga(u_\th, u_\th)| \le |\n_\th u_\th| + |\Ga(u_\th, \n_\th u_\th)| \le \eta + C_N \eta^2 \leq 2 \eta,\]
assuming that $\eta$ is sufficiently small (depending on $N$).
Thus
\[\begin{aligned}
  |\n_{\th} u_{\th}|^2&=|\p_\th u_\th +\Ga(u_\th, u_\th)|^2\\
  &\ge   |\p_\th u_\th|^2-2|\p_\th u_\th||\Ga(u_\th, u_\th)|+|\Ga(u_\th, u_\th)|^2\\
  &\ge  |\p_\th u_\th|^2-C \eta |u_\th|^2.
\end{aligned}\]
Then the ordinary Poincar\'e inequality on $S^1$ yields
\begin{equation}\label{f1overf_0}
f_1^2 =  \int_{S^1}|\n_{\th} u_{\th}|^2\ge  \int_{S^1}|u_\th|^2-C \eta \int_{S^1}|u_\th|^2=(1-C\eta )f_0^2.
\end{equation}
A similar argument, applied to $\nabla_\theta^2 u_\theta,$ gives
\begin{equation}\label{nablaf1overf1}
\int_{S^1}|\n_{\th}^2 u_{\th}|^2 \ge  (1-C\eta) f_1^2.
\end{equation}
\begin{comment}
Finally, by the Sobolev inequality on $S^1,$ we have
\begin{equation*}
\sup_{S^1} | u_\theta | \leq C \left( f_0 + f_1 \right) \leq C f_1.
\end{equation*}
\end{comment}
We first apply (\ref{nablaf1overf1}) and H\"older's inequality to (\ref{firstevol}), to obtain
\begin{equation}\label{secondevol}
\begin{split}
\frac12 \left( e^{2s} \partial_t f_1^2 - \p_s^2f_1^2 + (1 - C\eta) f_1^2 \right) & \leq - \int_{S^1} |\n_s \nabla_\theta u_{\th}|^2 - \int_{S^1} \LA I + II, \nabla_\theta u_\theta \RA \\
& \leq - \int_{S^1} |\n_s \nabla_\theta u_{\th}|^2 + \left( \|I\|_{L^2(S^1)} + \|II\|_{L^2(S^1)} \right) f_1.
\end{split}
\end{equation}
Note that
\[ \frac12\p_s f_1^2= f_1\p_s f_1=\int_{S^1}\<\n_s \n_\th u_{\th}, \n_\th u_\th\>\le \(\int_{S^1}|\n_s \n_\th u_{\th}|^2\)^{1/2}f_1,\]
so we have
$$|\p_s f_1|^2\le\int |\n_s \n_\th u_{\th}|^2.$$
On the other hand,
\[ \frac12\p_s^2(f_1^2)=f_1\cdot \p_s^2 f_1+ |\p_s f_1|^2.\]
Hence, ``dividing out'' by $f_1$ in (\ref{secondevol}) (which is justified in the distribution sense), we obtain
\begin{equation}\label{thirdevol}
\begin{split}
e^{2s} \partial_t f_1 - \p_s^2f_1 + (1 - C\eta) f_1 & \leq 2 \left( \|I\|_{L^2(S^1)} + \|II\|_{L^2(S^1)} \right).
\end{split}
\end{equation}
It remains to estimate the terms on the RHS of (\ref{thirdevol}). By (\ref{evolutions:derivbounds}), we have
\begin{equation*}%\label{smallness2}
e^{2s} \left(|u_t| + |\nabla_\theta u_t| + |\nabla^2_\theta u_t| \right) \leq C \eta.
\end{equation*}
Combining this with (\ref{smallness1}), from (\ref{I}), we obtain
\begin{equation*}
\left\| I \right\|_{L^2(S^1) } \leq C \xi_0 \eta e^{2s} + \eta^2 f_0 \leq C \xi_0 \eta e^{2s} + \eta f_1,
\end{equation*}
where we have used (\ref{f1overf_0}).
From (\ref{II}), since each term has at least one factor of $u_\theta$ or $\n_\theta u_\th,$ we also obtain
\begin{equation*}
\left\| I\! I \right\|_{L^2(S^1) } \leq C_N \eta^2 f_1 \leq \eta f_1,
\end{equation*}
for $\eta$ sufficiently small.
Inserting these estimates into (\ref{thirdevol}), and absorbing the $C \eta f_1$ terms in the LHS, we obtain
\begin{equation}\label{fourthevol}
\begin{split}
e^{2s} \partial_t f_1 - \p_s^2f_1 + (1 - C\eta) f_1 & \leq C \xi_0 \eta e^{2s}.
\end{split}
\end{equation}
Translating the above equation back to polar coordinates, we get (\ref{evolutions:angular}).






To estimate the radial energy
$$g = \sqrt{ \int_{ \{e^s\} \times S^1 } |u_s|^2 d\th },$$
we start from the identity
\begin{equation}\label{urevol:1}
 \frac12 \p_s^2 g^2= \int_{S^1}\( |\n_s u_s|^2+\< \n_s^2u_{s}, u_s\>\).
\end{equation}
Applying $\n_s$ to (\ref{T0ut}), we obtain
\begin{equation*}
\begin{split}
\nabla^2_s u_s + \nabla_s \nabla_\theta u_\theta & = \left( 2 \xi \partial_s \xi + 2 \xi^2 \right) e^{2s}  u_t + \xi^2 e^{2s} \nabla_s u_t \\
& = 2 \left( \xi^{-1} \partial_s \xi + 1 \right) \T_0(u) + \xi^2 e^{2s} \nabla_t u_s \\
& = 2 \left( \nabla_s u_s + \nabla_\theta u_\theta \right) + e^{2s} \nabla_t u_s + 2 \xi^{-1} \partial_s \xi (\nabla_s u_s + \nabla_\theta u_\theta)+ \left( \xi^2 - 1 \right) e^{2s} \nabla_s u_t.
\end{split}
\end{equation*}
We also have
$$\nabla_s \nabla_\theta u_\theta = \nabla_\theta^2 u_s + R(u_s, u_\theta) u_\theta.$$
Returning to (\ref{urevol:1}), we have
\begin{equation*}
\begin{split}
\frac12 \partial_s^2 g^2 & = \int |\nabla_s u_s|^2 + \int \LA e^{2s} \nabla_t u_s + 2 \nabla_s u_s + 2 \nabla_\theta u_\theta , u_s \RA \\
& \qquad + \int |\nabla_\theta u_s |^2 + \int \LA 2 \xi^{-1} \partial_s \xi \left( \nabla_s u_s + \nabla_\theta u_\theta \right) + (\xi^2 - 1) e^{2s} \nabla_s u_t , u_s \RA - \int R(u_s, u_\theta) u_\theta,
\end{split}
\end{equation*}
where we have integrated by parts once.
%We have
%$$ \int  \left( \nabla_s \log \xi + 1 \right) \LA \nabla_\theta u_\theta, u_s \RA = - 2 \int \LA \nabla_\theta \nabla_s \log \xi \, u_\theta, u_s \RA - 2 \int \left( \nabla_s \log \xi + 1 \right) \LA u_\theta, \nabla_\theta u_s \RA.$$
Applying H\"older's inequality, rearranging, and using (\ref{smallness1}), we obtain
\begin{equation*}
\begin{split}
\frac12 \left( e^{2s} \partial_t - \partial_s^2 + 2 \partial_s \right) g^2 & \leq - \int |\nabla_s u_s|^2  + 2 f_1 g + C \xi_0  \eta e^{2s} g \\
& \qquad - \int |\nabla_\theta u_s |^2 + C_N \eta f_0 g.
\end{split}
\end{equation*}
We choose $\eta$ small enough that $C_N \eta \leq 1,$ and discard the $-\int |\nabla_\theta u_s|^2$ term. After ``dividing out'' by $g$ as above, we obtain
\begin{equation}\label{urevol:2}
\begin{split}
\left( e^{2s} \partial_t - \partial_s^2 + 2 \partial_s \right) g \leq 6f_1 + C \xi_0 \eta e^{2s}.
\end{split}
\end{equation}
Changing back to polar coordinates and dividing by $r,$ we get the desired evolution equation.





\section{Radial heat kernel}




In this appendix, we extract several results from the appendix of \cite{lte}, replacing the integer dimension by a real number $\mu > 1.$ The proofs of Propositions \ref{festimate}-\ref{gestimate} are based on these results.

Let
$$\Delta_\mu = \p_r^2 + \frac{\mu - 1}{r} \p_r.$$
%n this appendix we port over some results from \cite{}, Appendix A, for the case of non-integer $\mu.$
%We write $H^\mu = H^\mu_{(0, \infty)}$ for 
%\begin{lemma} The function
%\begin{equation}
%H^\mu(r,s,t) = \frac{c_\mu e^{\frac{-(r^2 + s^2)}{4t}}}{t^{\mu/2} } \varphi\left( \frac{rs}{2t} \right),
%\end{equation}
%\end{lemma}
%\begin{proof}
%Supposing that $\mu$ is an integer, %we can obtain a kernel for the operator $\p_t - \Delta_\mu$ on $\left(0, \infty \right)$
In the case that $\mu$ is an integer, the spherical average of the Euclidean heat kernel is given by
\begin{equation}\label{Halphadef}
H(r,s,t) = \frac{c_\mu e^{\frac{-(r^2 + s^2)}{4t}}}{t^{\mu/2} } I\left( \frac{rs}{2t} \right),
\end{equation}
where $c_\mu$ is an appropriate constant, and\footnote{Since $I$ satisfies the ODE (\ref{Iode}), we in fact have
$$I(x) = x^{1 - \frac{\mu}{2}} I_{\frac{\mu}{2} - 1}(x),$$
where $I_{\frac{\mu}{2} -1}$ is the modified Bessel function. This recovers formula (3.3) of Bragg \cite{braggradialheatpolys}.}
\begin{equation}\label{varphidef}
I(x) = \int_0^\pi e^{x \cos \theta } \sin^{\mu - 2} \theta \, d\theta.
\end{equation}


\begin{lemma}\label{lemma:euclheatker} For any real $\mu > 1,$ the above function $H$ satisfies
\begin{equation*}
\begin{split}
\left( \p_t - \Delta_\mu \right) H (\cdot, s,t) = 0, \qquad t > 0, \\
H (r,s,t) > 0 \mbox{ for } 0 < r,s,t < \infty, \\
H (r, s,t) \to \frac{1}{s^{\mu - 1}} \delta(r-s) \qquad t \searrow 0,
\end{split}
\end{equation*}
and
\begin{equation}\label{Hbound}
\frac{C_\mu^{-1} e^{\frac{-(r-s)^2}{4t}}}{t^{1/2} (rs + t)^{\frac{n-1}{2}} } \leq H(r,s,t) \leq \frac{C_\mu e^{\frac{-(r-s)^2}{4t}}}{t^{1/2} (rs + t)^{\frac{n-1}{2}} }.
\end{equation}
\end{lemma}
\begin{proof} %For integer $\mu,$ this follows by construction.
%To check the formula for positive real $\mu,$ for a function of the form (\ref{Halphadef}),
We calculate
\begin{equation*}
\left( \p_t - \Delta_\mu \right) H(r,s,t) = \frac{c_\mu e^{\frac{-(r^2 + s^2)}{4t}}}{t^{\mu/2} } \left( \frac{-s^2}{4t^2} \right) \left( I''(x) + \frac{(\mu - 1)}{x} I'(x) - I(x) \right).
\end{equation*}
For $I(x)$ given by (\ref{varphidef}), we have
\begin{equation}\label{Iode}
\begin{split}
I''(x) + \frac{(\mu-1)}{x} I'(x) - I(x) & = \int_0^\pi e^{x \cos \theta} \left( \cos^2 \theta + \frac{(\mu-1)}{x}\cos\theta - 1 \right) \sin^{\mu - 2} \theta \, d\theta \\
& = \int_0^\pi e^{x \cos \theta} \left( -\sin^2 \theta + \frac{(\mu - 1)}{x}\cos\theta\right) \sin^{\mu - 2} \theta \, d\theta \\
& = 0,
\end{split}
\end{equation}
after integrating by parts.
%$$\varphi''(x) + \frac{n-1}{x} \varphi'(x) - \varphi(x) = 0,$$
Hence, $H$ solves the PDE as required for any real $\mu.$

Borrowing a factor of $e^{\frac{rs}{2t}}$ in (\ref{varphidef}), we have
\begin{equation}\label{HI1}
H(r,s,t) = \frac{c_\mu e^{\frac{-(r - s)^2}{4t}}}{t^{\mu/2} } I_1\left( \frac{rs}{2t} \right),
\end{equation}
where
\begin{equation*}%\label{varphidef}
I_1(x) = \int_0^\pi e^{x (\cos \theta - 1) } \sin^{\mu - 2} \theta d\theta.
\end{equation*}
Then $I_1(x)$ clearly tends to a positive constant as $x \to 0.$ By the substitution $u = \sqrt{x(1 - \cos \theta )},$ it follows that the integral is bounded by a constant times $x^{-\frac{\mu - 1}{2}}.$ Hence
$$I_1(x) \leq \frac{C}{(1 + x)^{\frac{\mu - 1}{2}}} \leq \left( \frac{t}{rs + t} \right)^{\frac{\mu - 1}{2}}.$$
Substituting into (\ref{HI1}), we obtain the desired bound.
\end{proof}



\subsection{Initial data}\label{app:initialkernel} Let $H_{\LB \rho, R \RB}(r,s,t)$ be the Dirichlet kernel for the operator $\p_t - \Delta_\mu$ on the interval $\LB \rho, R\RB,$ satisfying
\begin{align*}
\left( \p_t - \Delta_\mu \right) H_{\LB \rho, R\RB}(\cdot, s,t) & = 0 & t > 0, \\
H_{\LB \rho, R\RB}(\rho, s,t) & = 0 = H_{\LB \rho, R\RB} (R, s,t) &  \rho \leq s \leq R , \,\, t > 0, \\
H_{\LB \rho, R\RB}(r, s,t) & \to \frac{1}{s^{\mu - 1}} \delta(r - s) &  t \searrow 0.
\end{align*}
By the maximum principle, we have $0 \leq H_{\LB \rho, R \RB}(r,s,t) \leq H(r,s,t),$ and so
\begin{equation}\label{HrhoRbound}
0 \leq H_{\LB \rho , R \RB} (r,s,t) \leq \frac{C_\mu e^{\frac{-(r-s)^2}{4t}}}{t^{1/2} (rs + t)^{\frac{\mu-1}{2}} }.
\end{equation}
Given an initial function $\varphi(r)$ on $\LB \rho, R \RB,$ the solution of the initial-value problem is given by
\begin{equation}\label{ivp}
v_0(r,t) = \int_\rho^R H_{\LB \rho , R \RB} (r,s,t) \varphi(s) s^{\mu - 1} \, ds.
\end{equation}
Let
$$w^a(r,t) = \left( \frac{r^2}{r^2 + t} \right)^{a/2}.$$

\begin{prop}\label{initialestimate} For $0 \leq k \leq \mu - 1,$ assuming that $|\varphi(r) | \leq A r^{-k},$ we have
\begin{equation*}
|v_0(r,t)| \leq C_\mu A r^{-k} w^{k}(r,t) w^{\mu - k}(R,t).
\end{equation*}
\end{prop}
\begin{proof} From (\ref{HrhoRbound}) and (\ref{ivp}), we have
\begin{equation*}
\begin{split}
|v_0(r,t)| & \leq C A \int_\rho^R e^{-(r-s)^2/4t} \frac{s^{-k + \mu -1}}{(rs + t)^{\frac{\mu -1}{2}}} \, \frac{ds}{t^{1/2}} \\
& \leq CA r^{-k} \int_\rho^R e^{-(r-s)^2/4t} \frac{r^{k} s^{-k + \mu - 1}}{(rs + t)^{\frac{\mu -1}{2}}} \, \frac{ds}{t^{1/2}}.
\end{split}
\end{equation*}
By Lemma A.1a of \cite{lte}, applied with $a = k, b= \mu - k -1,$ and $c = d = 0,$ we have
\begin{equation*}
\begin{split}
\int_\rho^R e^{-(r-s)^2/4t} \frac{r^{k} s^{\mu - k - 1}}{(rs + t)^{\frac{\mu -1}{2}}} \, \frac{ds}{t^{1/2}} & \leq C \frac{ R - \rho}{R - \rho + \sqrt{t} } w^{k}(r,t) w^{\mu - k - 1}(R,t) \\
& \leq C w^{k}(r,t) w^{\mu - k}(R,t).
\end{split}
\end{equation*}
The result follows.
\end{proof}



%With the bound of Lemma \ref{lemma:heatbound} in hand, we can carry over the estimates of \cite{}, Appendix A, with minimal changes.

\subsection{Boundary data}\label{app:innerkernel} To construct a kernel for the boundary data at the inner radius $\rho = 1,$ we follow the argument of \cite{lte}, Appendix A.3. Suppose $R > 1,$ and let
$$h(r)  = \frac{r^{2 - \mu} - R^{2 - \mu} }{1 - R^{2 - \mu} }.$$
Let
$$y_1(r,t) = h(r) - \int_1^R H_{\LB 1, R \RB}(r,s,t) h(s) \, s^{\mu - 1} ds.$$
This satisfies
\begin{equation*}
\begin{split}
(\p_t - \Delta_\mu) y_1 & = 0 \\
y_1(r,0) & = 0, \quad 1 < r < R \\
y_1(1,t) & = 1, \,  y_1(R,t) = 0, \quad t > 0.
\end{split}
\end{equation*}
The function
$$G_{\LB 1, R\RB}(r,t) = \p_t y_1(r,t)$$
satisfies
$$\lim_{r \searrow 1} G_{\LB 1, R\RB}(r,t) = \delta(t).$$
\begin{lemma}\label{boundarykernel}
We have
\begin{equation*}
0 \leq G_{\LB 1, R\RB}(r,t) \leq \frac{C_{\mu} e^{\frac{-(r -1)^2}{5t} }}{t(t + 1)^{\frac{\mu}{2} - 1}} \cdot \begin{cases} \min \LB (r-1)/\sqrt{t}, 1 \RB & (t \leq 1) \\
\min \LB r - 1, 1 \RB & (t \geq 1). \end{cases}
\end{equation*}
\end{lemma}
\begin{proof} Replacing $n$ by $\mu,$ the bound is identical to that of Lemma A.4b of \cite{lte}, and the proof there carries over. 
\end{proof}

To obtain an inner boundary kernel for $\LB \rho, R \RB,$ we let
$$G_{\LB \rho, R \RB}(r,t) = \frac{1}{\rho^{2}} G_{\LB 1, R/\rho\RB}(r/\rho,t/\rho^2).$$
By Lemma \ref{boundarykernel}, this satisfies
\begin{equation}\label{Gbound}
G_{\LB \rho, R \RB}(r,t) \leq \frac{C_{\mu} e^{\frac{-(r -\rho)^2}{5t} } \rho^{\mu - 2} }{t(t + \rho^2 )^{\frac{\mu}{2} - 1}} \cdot \begin{cases} \min \LB (r-\rho)/\sqrt{t}, 1 \RB & (t \leq 1) \\
\min \LB r - \rho, 1 \RB & (t \geq 1). \end{cases}
\end{equation}
The solution of the boundary problem with data $\psi(t)$ at $r = \rho$ is given by
\begin{equation}\label{boundaryivp}
v_1(r,t) = \int_0^t \psi(\tau) G_{\LB \rho, R \RB}(r,t - \tau) \, d\tau.
\end{equation}

\begin{prop}\label{innerestimate} For $2\rho \leq r \leq R$ and $t \geq 0,$ we have
\begin{equation*}
| v_1(r,t) | \leq C_\mu e^{-(r - \rho)^2/6t} \left( \frac{\rho^{ \mu - 2 }}{r^{\mu - 1}} \right) \sqrt{ \int_0^t \psi^2(\tau) d\tau }.
\end{equation*}
\end{prop}
\begin{proof} We apply H\"older's inequality as in the proof of Proposition A.5b of \cite{lte}. From (\ref{Gbound}-\ref{boundaryivp}), we have:
\begin{equation*}
\begin{split}
r^{\mu - 1} |v_1(r,t)| & \leq C_\mu \rho^{ \mu  - 2} \int_0^t |\psi(\tau)| \frac{ r^{\mu - 1} e^{\frac{-(r -\rho)^2}{5(t - \tau)} }}{(t - \tau)(t - \tau + \rho^2)^{\frac{\mu}{2} - 1}} \, d\tau \\
& \leq C_\mu \rho^{ \mu - 2} \sqrt{\int_0^t \psi^2(\tau) \, d\tau } \\
& \qquad \cdot \sqrt{ \int_0^t \frac{r^{2\mu-2}(r - \rho)^2 }{(t - \tau)^2(t - \tau + \rho^2)^{\mu - 2} } e^{\frac{-2(r -\rho)^2}{5(t - \tau)} } \, \frac{d\tau}{(r - \rho)^2} }.
\end{split}
\end{equation*}
The result follows by changing variables $u = \frac{\tau}{(r  - \rho)^2}.$ % one sees that the second integral is bounded by a constant times $e^{-(r - \rho)^2/6t}.$ The result follows.
\end{proof}






\bibliographystyle{amsinitial}
\bibliography{biblio}

\end{document}







\begin{thebibliography}{20}

\bibitem{AHM2009}
S. B. Angenent, J. Hulshof, and H. Matano, {\em The radius of vanishing bubbles in equivariant harmonic map flow from $D^2$ to $S^2$}. SIAM J. Math. Anal. 41, no. 3, 1121-1137 (2009).


\bibitem{BairdEells1964}
P. Baird and J. Eells. \emph{A conservation law for harmonic maps.} Geometry Symposium Utrecht 1980, Springer (1981).

\bibitem{BHK2003}
J. B. van den Berg, J. Hulshof, and J. R. King, {\em Formal asymptotics of bubbling in the harmonic map heat flow.} SIAM J. Appl. Math., 63, pp. 1682-1717 (2003).

%\bibitem{ChenSong2020}
%B. Chen, C. Song, {\em Isolated singularities of Yang-Mills-Higgs fields on surfaces}, Inter. Math. Res. Not. (2020) arXiv:1907.07092.

%\bibitem{ChenLi2014}
%J. Chen, Y. Li, {\em Homotopy classes of harmonic maps of the stratified 2-spheres and applications to geometric flows}. Adv. Math. 263 (2014), 357-388.

\bibitem{CDY1992}
K. Chang, W. Ding and R. Ye, \emph{Finite-time blow up of the heat flow of harmonic maps from surfaces}. J. Diff. Geom. \textbf{36}, 507-515 (1992).

\bibitem{DPW2020}
J. D\'{a}vila, M. del Pino, and J. Wei, \emph{Singularity formation for the two-dimensional harmonic map flow into {$S^2$}}. Invent. Math. 219, no. 2, 345-466 (2020).

\bibitem{DingTian1995}
W. Ding and G. Tian, \emph{Energy identity for a class of approximate harmonic maps from surfaces}. Comm. Anal. Geom. \textbf{3}, 543-554 (1995).

\bibitem{EellsSampson1964}
J. Eells and J. H. Sampson, \emph{Harmonic mappings of Riemannian manifolds.} Amer. J. Math. 86, 109-169 (1964).

\bibitem{GoldbergKobayashi1967}
S. I. Goldberg, and S. Kobayashi, \emph{Holomorphic bisectional curvature.} J. Differ. Geom. 1.3-4: 225-233 (1967).


%\bibitem{Gu2009}
%H. Gu, \emph{A new proof of Mok's generalized Frankel conjecture theorem}. Proc Amer Math Soc, 2009, 137: 1063-1068

%\bibitem{HongYin2013}
%M.-C. Hong and H. Yin, \emph{On the Sacks-Uhlenbeck flow of Riemannian surfaces}, Comm. Anal. Geom. 21 (2013), no. 5, 917-955.

\bibitem{LinWang1998}
F. Lin and C. Wang, \emph{Energy identity of harmonic map flows from surfaces at finite singular time}. Calc. Var. \textbf{6}, 369-380 (1998).

\bibitem{LiuYang2010}
Q. Liu and Y. Yang, {\em Rigidity of the harmonic map heat flow from the sphere to compact K\"ahler manifolds}. Ark. Mat. 48, no. 1, 121-130 (2010).

\bibitem{Mok1988}
N. Mok, \emph{The uniformization theorem for compact K\"{a}hler manifolds of nonnegative bisectional curvature}. J Diff Geom, 27: 179-214 (1988).

\bibitem{Mori1979}
S. Mori, \emph{Projective manifolds with ample tangent bundles}. Ann of Math, 110: 593-606 (1979).

\bibitem{Parker1996}
T. Parker, \emph{Bubble tree convergence for harmonic maps}. J. Diff. Geom. \textbf{44}(3), 595-633 (1996).

\bibitem{Parker2003}
T. H. Parker. \emph{What is a Bubble Tree?} Notices of the American Mathematical Society, 50(6), 666-667 (2003).

\bibitem{QingTian1997}
J. Qing and G. Tian, \emph{Bubbling of the heat flows for harmonic maps from surfaces}. Commun. Pure Appl. Math. 50(4), 295-310 (1997).

%\bibitem{Rade1992}
%J. R{\aa}de, \emph{Decay estimates for Yang-Mills fields; two new proofs}. Global analysis in modern mathematics (Orono, ME, 1991; Waltham, MA, 1992), 91-105, Publish or Perish, Houston, TX, 1993.

\bibitem{RaphaelSchweyer2013}
P. Rapha\"{e}l and R. Schweyer, {\em Stable blowup dynamics for the 1-corotational energy critical harmonic heat flow}. Comm. Pure Appl. Math. 66, no. 3, 414-480 (2013).

\bibitem{RaphaelSchweyer2014}
P. Rapha\"{e}l and R. Schweyer, {\em Quantized slow blow-up dynamics for the corotational energy-critical harmonic heat flow}. Anal. PDE 7, no. 8, 1713-1805 (2014).

\bibitem{SacksUhlenbeck1981}
J. Sacks and K. Uhlenbeck, \emph{The existence of minimal immersions of 2-spheres}. Ann. of Math. 113 1-24 (1981).

\bibitem{SchoenUhlenbeck1982}
R. Schoen and K. Uhlenbeck. \emph{A regularity theory for harmonic maps.} J. Differ. Geom. 17.2: 307-335 (1982).

\bibitem{SchoenYau1978}
R. Schoen and S. T. Yau, \emph{On univalent harmonic maps between surfaces.} Invent. Math., 44(3), 265-278 (1978).

\bibitem{SchoenYau1997}
R. Schoen and S. T. Yau, {\em Lectures on harmonic maps. Conference Proceedings and Lecture Notes in Geometry and Topology, II. } International Press, Cambridge, MA, (1997).

\bibitem{SiuYau1980}
Y. T. Siu and S. T. Yau, \emph{Compact K\"{a}hler manifolds of positive bisectional curvature}. Invent. Math. 59: 189-204 (1980).

\bibitem{Struwe1985}
M. Struwe, \emph{On the evolution of harmonic mappings of Riemannian surfaces}. Comment. Math. Helv. 60(4), 558-581 (1985).

\bibitem{struwevariationalmethods}
M. Struwe, \emph{Variational Methods, 3rd ed. Applications to Nonlinear Partial Differential Equations and Hamiltonian Systems}. Berlin: Springer-Verlag (1996).

\bibitem{Toledo1979}
D. Toledo, {\em Harmonic maps from surfaces to certain Kaehler manifolds.} Math. Scand. 45, no. 1, 13-26 (1979).

\bibitem{Topping1997}
P. Topping, {\em Rigidity in the harmonic map heat flow.} J. Differential Geom. 45 no. 3, 593-610 (1997).

\bibitem{Topping2004-mathz}
P. Topping, {\em Winding behaviour of finite-time singularities of the harmonic map heat flow.} Math. Z. 247, no. 2, 279-302 (2004).

%\bibitem{Topping2004-annmath}
%P. Topping, {\em Repulsion and quantization in almost-harmonic maps, and asymptotics of the harmonic map flow}, Ann. of Math. (2) 159 (2004), no. 2, 465-534.

\bibitem{Topping2004-cvpde}
P. Topping, {\em Improved regularity of harmonic map flows with H\"{o}lder continuous energy}. Calc. Var. Partial Differential Equations 21, no. 1, 47-55 (2004).

\bibitem{Waldron2016}
A. Waldron, {\em Instantons and singularities in the Yang-Mills flow.} Calc. Var. Partial Differential Equations 55 (5), 113 (2016).

\bibitem{Waldron2019}
A. Waldron, {\em Long time existence for Yang-Mills flow.} Invent. Math. 217, no. 3, 1069-1147 (2019).

\bibitem{Wang1996}
C. Wang, \emph{Bubble phenomena of certain Palais-Smale sequences from surfaces to general targets.} Houston. J. Math. 22(3), 559-590 (1996).

\bibitem{Wood1979}
J. C. Wood, \emph{Holomorphicity of certain harmonic maps from a surface to complex projective n-space.} Journal of the London Mathematical Society 2.1: 137-142 (1979).


